\section[Notebook 1968-13: the space $X \cup_\Phi \{\infty\}$ and the smash lemma]{Notebook 1968-13: the space $X \cup_\Phi \{\infty\}$ of a family of supports, and the smash lemma\protect\piste\protect\verifie{Notebook1968_13}}\label{sec:quillen-1968-13}

\emph{Context.} In December 1968, \carnet{1968}{13} relates supports and bordism. On page~64, Quillen takes on a space $X$ a family of supports $\Phi$: a family of closed sets, stable under finite unions and under passage to closed subsets. On page~65, he forms the space $X \cup_\Phi \{\infty\}$, in which $X$ is open and the neighbourhoods of $\infty$ are the sets $(X \setminus Z) \cup \{\infty\}$, $Z \in \Phi$. For $\Phi$ the family of closed sets, $\infty$ is an isolated point; for $\Phi$ the family of compact sets, one recovers the Alexandroff compactification $X^+$. A crossed-out passage on page~64 notes that the space is Hausdorff if $\Phi$ contains a neighbourhood of each point. On pages~67--69, after his own correction (\q{Correction}: he had defined the smash product wrongly, it being the quotient of $X \times Y$ by the wedge), he states and proves that $X^+ \wedge (Y \cup_\Phi \{\infty\}) = (X \times Y) \cup_\Psi \{\infty\}$, where $\Psi$ is the family of closed subsets of $X \times Y$ contained in some $X \times F$, $F \in \Phi$, and proper over $F$.

\emph{Intuition.} This is the argument of the page. A neighbourhood of $\infty$ in the smash product is the image of an open set containing the wedge $X^+ \times \{\infty\} \cup \{\infty\} \times (Y \cup_\Phi \{\infty\})$; its complement is a closed set $C$ of $X \times Y$, closed in the whole product. Since $X^+$ is compact, the tube lemma keeps $C$ away from $X^+ \times \{\infty\}$, that is, places it inside some $X \times F$; and $C$ is closed in $X^+ \times F$ (figure~\ref{fig:quillen-1968-13}).


\begin{figure}[htbp]
\centering
\begin{tikzpicture}[every node/.style={font=\small}]
\fill[black!8] (0,0.7) rectangle (6,2.3);
\draw[black!30] (0,0.7) -- (6,0.7) (0,2.3) -- (6,2.3);
\draw[black!55] (0,0) -- (6,0) (0,0) -- (0,4.2);
\draw[line width=1.8pt] (0,4.2) -- (6,4.2) -- (6,0);
\fill[frise!20] (1.0,1.3) .. controls (1.4,0.8) and (3.2,0.8) .. (4.2,1.2) .. controls (4.9,1.5) and (4.6,2.1) .. (3.6,2.0) .. controls (2.6,1.9) and (1.6,2.15) .. (1.2,1.9) .. controls (0.9,1.7) and (0.85,1.5) .. cycle;
\draw[frise] (1.0,1.3) .. controls (1.4,0.8) and (3.2,0.8) .. (4.2,1.2) .. controls (4.9,1.5) and (4.6,2.1) .. (3.6,2.0) .. controls (2.6,1.9) and (1.6,2.15) .. (1.2,1.9) .. controls (0.9,1.7) and (0.85,1.5) .. cycle;
\node[frise] at (2.8,1.45) {$C$};
\draw[decorate, decoration={brace, amplitude=4pt}] (-0.12,0.7) -- (-0.12,2.3) node[midway, left=5pt] {$F$};
\fill (6,0) circle (1.6pt) node[below] {$\infty$};
\fill (0,4.2) circle (1.6pt) node[left] {$\infty$};
\node[below] at (3,-0.05) {$X$};
\node[left] at (-0.05,3.3) {$Y$};
\node[anchor=west] at (6.25,2.6) {$\{\infty\} \times (Y \cup_\Phi \{\infty\})$};
\node[above] at (3,4.25) {$X^+ \times \{\infty\}$};
\node[black!60] at (5.25,1.0) {$X \times F$};
\end{tikzpicture}
\caption{The product $X^+ \times (Y \cup_\Phi \{\infty\})$, with $X^+$ on the horizontal axis and $Y \cup_\Phi \{\infty\}$ on the vertical axis. The wedge, in thick line, is the union of the upper edge $X^+ \times \{\infty\}$ and the right edge $\{\infty\} \times (Y \cup_\Phi \{\infty\})$. The complement of a neighbourhood of $\infty$ in the smash product is a closed set $C$ of $X \times Y$ (in blue). The tube lemma, applied to the compact set $X^+ \times \{\infty\}$, places it in a strip $X \times F$ with $F \in \Phi$ (in grey), and $C$ is closed in $X^+ \times F$: it does not accumulate on the right edge.}
\label{fig:quillen-1968-13}
\end{figure}

\begin{encadre}
\begin{theoreme}[Quillen, notebook 1968-13, pp.~64--69]\label{thm:quillen-1968-13}
Let $X$, $Y$ be topological spaces and $\Phi$ a family of supports on $Y$.
\begin{enumerate}
\item[a)] $X$ is an open subset of $X \cup_\Phi \{\infty\}$ and $\{\infty\}$ is closed. If $X$ is Hausdorff, $X \cup_\Phi \{\infty\}$ is Hausdorff if and only if every point of $X$ has a neighbourhood in $\Phi$.
\item[b)] For $\Phi$ the family of compact closed sets, $X \cup_\Phi \{\infty\}$ is homeomorphic to $X^+$; for $\Phi$ the family of closed sets, $\infty$ is isolated.
\item[c)] Let $\Psi$ be the family of closed subsets $A$ of $X \times Y$ contained in some $X \times F$, $F \in \Phi$, and closed in $X^+ \times Y$. The map $X^+ \times (Y \cup_\Phi \{\infty\}) \to (X \times Y) \cup_\Psi \{\infty\}$, which is the identity on $X \times Y$ and sends the wedge to $\infty$, is a quotient map: $(X \times Y) \cup_\Psi \{\infty\}$ is the smash product $X^+ \wedge (Y \cup_\Phi \{\infty\})$ with the quotient topology.
\end{enumerate}
\end{theoreme}
\end{encadre}

\noindent Quillen assumes $X$ locally compact, which makes $X^+$ Hausdorff; the proof uses only the compactness of $X^+$, and c) holds without any hypothesis. The family $\Psi$ is taken in the form that the proof on the page uses (closed in $X^+ \times F$); the equivalence with \q{proper over $F$}, which requires $X^+$ Hausdorff, is not formalized.

\begin{proof}
a) Separating $x \in X$ from $\infty$ requires an open set $U \ni x$ disjoint from some $X \setminus Z$, $Z \in \Phi$, that is, $U \subset Z$: $Z$ is a neighbourhood of $x$. Conversely, the interior of $Z$ and $(X \setminus Z) \cup \{\infty\}$ separate $x$ and $\infty$.

b) An open set containing $\infty$ has, in both topologies, a complement in $X$ that is closed and compact, or contained in a compact closed set, which amounts to the same thing. For the closed sets, $X$ itself is in $\Phi$, and $\{\infty\}$ is open.

c) The map is surjective. For a set $U$ not containing $\infty$, its preimage is the image of $U \cap (X \times Y)$ under the open embedding $X \times Y \to X^+ \times (Y \cup_\Phi \{\infty\})$. For $U \ni \infty$, the preimage is the complement of the image of the closed set $C = (X \times Y) \setminus U$, and everything comes down to seeing that this image is closed if and only if $C \in \Psi$. If it is closed, the tube lemma, applied to the compact set $X^+ \times \{\infty\}$, gives $F \in \Phi$ with $C \subset X \times F$, and the preimage under $\mathrm{id} \times \iota : X^+ \times Y \to X^+ \times (Y \cup_\Phi \{\infty\})$ shows that $C$ is closed in $X^+ \times Y$. Conversely, the image of $C$ is contained in the closed set $X^+ \times F$, which does not meet $\{\infty\}$, and it is there the image of a closed set under an embedding.
\end{proof}

\begin{sloppypar}
\noindent\emph{Verification in Lean.} \texttt{SupportFamily} and \texttt{Extension} define the family and the space; a) is \texttt{isOpenEmbedding\_ι}, \texttt{isClosed\_infty} and \texttt{t2Space\_iff}; b) is \texttt{homeomorphOnePoint}, to mathlib's compactification \texttt{OnePoint X}, \texttt{isOpen\_infty\_closedSupports} and \texttt{homeomorphSum}, to the sum of $X$ and a point; c) is \texttt{isQuotientMap\_smashMap}, with the family \texttt{smashSupports} and the key lemma \texttt{isClosed\_image\_e\_iff}, in \href{\leanblobmain/Notebook1968_13.lean}{\texttt{Notebook1968\_13.\allowbreak{}lean}}.
\end{sloppypar}

Not formalized: $\Gamma(X^+, j_!F) = \Gamma_\Phi(X, F)$ (p.~65), the exact sequence that follows from it (p.~66), and the application to bordism (pp.~66--69).

\litterature{the space is found in the literature, the smash lemma is not. G.~Beer and M.~C.~Vipera associate to every family of ideals of subsets with a closed base of a space an extension of that space, of which the Alexandroff compactification is the prototype \cite{beervipera2012}; a family of supports generates such an ideal, and the one-point extension corresponding to it is $X \cup_\Phi \{\infty\}$. G.~Beer was already studying the one-point extension of a space equipped with a bornology \cite{beer2008}, a setting that goes back to S.-T.~Hu \cite{hu1949}. These references rest on the zbMATH Open reviews (Zbl~1259.54008 and 1171.54012); the full text has not been read, and it has not been checked that the separation criterion a) appears there in this form. Supports and $\Gamma_\Phi$ are due to Godement (1958) and to Bredon (\emph{Sheaf Theory}), whose book has not been combed page by page. The smash lemma has not been found in print.}
